\chapter{Doubt}

\vspace{-1.3cm}
\begin{localsize}{10}
	\begin{quote}
		``I have never made but one prayer to God, a very short one: Oh Lord, make my enemies ridiculous. And God granted it."
		\begin{flushright}- Voltaire \end{flushright}
	\end{quote}
\end{localsize}
\vspace{1cm}


``One day is today," Kraerion whispered to himself as his tired eyes slowly widened.

The chipmunk had woken---like all mornings---to a loud, unpleasant banging sound, echoing around the run outside the comforts of his small burrow.
The sound carried with it the news of a new day, and with it the message that all who wanted a day's pay ought to drag themselves to the great hall, the Aorta, presently.
The frowsty air this deep underground usually made him rise as soon as he woke to seek a more pleasant ambiance---but today he lingered in thought.

``One day is today," he repeated, to reassure himself and still his doubts, as if repeating the words would make them more true.
The power of inertia is strong; Kraerion knew full well of its insidious nature---that \textit{tomorrow} and \textit{never} were but equal measurements of eternity.
He had longed---even dreamed---for this day to come for as long as he could remember.
He wondered why, now when the day was finally here, he had so suddenly lost his strong convictions.

Chipmunks may be genetically poised for a life underground, but not the chipmunk Kraerion.
Even as a little kitten, he had loathed their existence: the damp runs and burrows, the secondhand air, the ever-present darkness, and just about everything else underground by mere association.
He hated---as only an adolescent can hate---that which most of his peers called home.
Instead, he wished for a strong wind in his fur, sunlight that uncovered all the beauty of Nature, but most of all, the chaotic buzzing of life above, that so contrasts with the routines down below.
He yearned to run the branches of trees, not the ever-branching network of runs dug in the forest soil.
His father had joked many a time that his son was born more a squirrel than chipmunk---and he had not been wrong.
When he was younger, he had declared that one day he would live amongst the trees, but to that received feedback usually only appropriate after a well-executed joke.

Kraerion eventually mustered the will to rise from his moss-coated bed.
He went over to a corner of his burrow and began scratching away at the dirt.

``Good, still here," he confirmed anxiously, as he felt the cold touch of seven steel screw-nuts---which would surely appear beautiful in the sunlight---lying before him.
He covered them up again, making sure that no one could find them.
The chipmunk left his burrow, stepped into the run, and followed its twists and turns until it widened into a great hall: the Aorta.
It was dug under a great willow; its massive network of roots---intertwined with the walls and ceiling---supported the weight of the dirt and kept it from caving in.\\

Since the time of Man, the forest had steadily increased its economic output.
One of the contributing factors was the increased specialization of its inhabitants.
Whilst many a species have some certain advantage in the gathering of some specific type of feed, none is as great a carrier as the chipmunk.
In their elastic cheeks, they are able to carry loads of feed---and other goods for that matter---efficiently around the forest.
Their small size and great numbers gave the operation a granularity to it, which made it difficult for larger and stronger animals like the roebuck to compete, except for very large packages.
When other animals began to hire them for their services, the chipmunks expanded their network of tunnels and runs to cover the whole forest, connecting all its major hubs.
Therefore, every morning, the chipmunks gathered in the Aorta to share news and plan the day ahead.

Most chipmunks were already present when Kraerion entered the hall.

``One day is today," he reminded himself again, as he headed towards his usual spot in the front.
But he hesitated suddenly and reconsidered; instead, he decided to begin his new life by breaking this habit, finding a place in the very back.\\

The murmur of conversation disappeared in an instant when an old chipmunk appeared in the Aorta.
His name was Aequitius, and he was the chief chipmunk.
Gray spots of fur spoke of his age, yet his features were easy on the eye, and he bore his years proudly.
Everyone knew the chief and saw him as something akin to a revered hero.
When he was but a kitten, he'd taken interest in the chipmunks' livelihood and noticed irregularities in their billings.
After some further deliberation with his best friend Br'er Rabbit, he realized that they could likely double their prices with little to no churn.
He put forth the motion to the then-chief chipmunk, who surprisingly listened to the youngster and approved it.
When their income subsequently increased, Aequitius was rather vocal about its occurrence, and made it impossible for the chief chipmunk not to pass it forward to the chipmunks' daily wage.

When Aequitius a few years later became unanimously chosen as the next chief, he took it a step further.
He transformed the chipmunks' enterprise to a sort of cooperative, with bylaws dictating that a fixed percentage of the total revenue would be earmarked for workers' wages.
He also instituted rules limiting the influence of his own office---not only for himself, but setting a precedent for his eventual replacements.

Together with Br'er Rabbit---who had by then become an informal leader of the rabbits---and with the help of the by then-Twig of the Glade of Representatives, Ongenþeow, he started the Winter Fund.
It was a collective fund that functioned similarly to human insurance: its members contributed feed during the spring, summer, and autumn as a premium; members who found themselves short during the winter could then make a claim and receive feed to sustain them until winter's end.
%Dearth

A cynic human might proclaim the flaw in this scheme---that an animal might take advantage of it by wasting his wage and making a claim every winter---but he forgets two things: most animals can't count; and although possible, making a claim is very stigmatized.
Most animals would rather starve than to broadcast their destitution.
Aequitius---being a co-founder of the fund---also made sure that all the work related to its logistics went to the chipmunks.
Moving and storing the feed deep underground---where the soil is grainy, cold, and preserves the feed from pest and rot whilst keeping it hidden from prying thieves---further boosted the demand for their services.

The fund was a great success; winter deaths decreased drastically the first year, and then slowly year after year.
The chipmunks secured a steady stream of sophisticated work from the fund, and thus means for a better life.
For all the contributions Aequitius made, none could match the feeling he instilled in every chipmunk: the intangible sense of worth, the standing it brought them in the forest, and the care he fostered between chipmunks by treating them all as equals.\\

When Aequitius reached the front of the Aorta, he turned and faced the scurry.

``Before we go through today's agenda, I have news---both good and bad," he began.
``And, as you may not be able to appreciate the good in anticipation of the bad, I'll begin with the bad."
He took a tactful pause, as anyone with his level of experience would in this situation.

``The eastbound run collapsed late last evening; we're completely cut off from the eastern part of the forest.
Therefore, the deliveries to the rabbits' hill must be postponed."

``Why did it collapse?" Aequitius's younger brother Pietus asked, who stood at his regular place just right of his brother.

``Well---" Aequitius sighed. ``A pheasant went off chasing a maggot again and crossed above the shallow no-zone."

``Again!?" Pietus cried, not hiding his contempt.
``Did the pheasants not promise to be more careful, that last time would be the last?
Will they honour our agreement and pay for the damages?"

``They did indeed.
As you all know, the pheasants are not the most organized of animals and will only blame each other if we make a claim for damages.
We may bring it to the Glade, but what---other than appearing petty---would we gain?
Little and less, I'm afraid.
Nevertheless, I want you, Pietus, to go to the moles and have them dig a deeper eastbound run this time.
That is, at least, a factor we can affect."
``I'll do it," Pietus responded.
``But send someone else to the moles next time.
They don't even speak our language..." %the common tongue.

``Nā, hēo seċġaþ þone ealda sprǣċe.
It is we that don't speak theirs." Aequitius paused, waiting to see if Pietus had a response.
When his brother only mumbled something inaudible, he continued.

``I'll make sure Br'er and the rabbits get the news of this postponement personally.
As for the good news, the rumors going around are true: the Glade has decided to institute a 2\% tax on all Forest shipments.
It's only a temporary experiment with us as the pilot service."

``How is this good news?" Veritas squeaked, a tiny runt of a chipmunk who was also standing in the front close to Aequitius.

``Two reasons: firstly, the funds are marked as aid for destitute animals who can't afford the Winter Fund, and we'll surely get the contract for its operation.
Secondly, we don't pay for our own services---we're essentially exempt!"

``No! That's not how it works." Veritas snapped angrily. ``This will force an increase in our prices or hurt our margins.
To only burden one industry, and ours at that... it's incredibly unfair."

``Enough," Aequitius grumbled.
``That's enough, this is neither the place nor time; we may discuss this later in private."

Veritas did not answer, but she turned her head sideways---showing everyone her disdain.

The old chipmunk moved on to the day's normal agenda, dishing out more routine orders before dismissing his workforce.
However, his eyes were fixed on Kraerion---the dismissal did not apply to him.
When the last chipmunk had left the Aorta, save for the two of them, he began to speak.

``Kraerion---" he said, now very serious.

``Father." Kraerion met the old chipmunk's eyes, matching his tone.

``It's true then, hiding in the back like that?
You're actually leaving what we worked so hard to create?"

``What \textit{you} created, you mean?
You cannot pretend this is news. I've been more than forthcoming with my intention to lead a different life.
You of all chipmunks should know of my extra shifts to afford this move, and my refusal of the managerial responsibilities you so blatantly pushed on me."

``We all have dreams, Kraerion. I know you might \textit{feel} like a squirrel, but you're not one.
%dearth of resources
I promise you that you'll end up miserable, bereft of resources, and completely alone in a world of wings if you try to live as one.
Dreams are a lesson in growing up: learning how to smother them in favour of what's truly important in this life. And, you'll get other dreams---realistic dreams---like forming a family, building a future for our species, and contributing to the betterment of the Forest at large."

``I'm sorry," Kraerion said at length, ``but that's not me.
And quite frankly, that'll never be me."

``Don't you feel any sense of responsibility? What about your sister, Veritas?
You expect Pietus and me to go on forever?"

``What about her? She's smarter than us three combined; she could take over tomorrow if you'd only let her."

``You know full well that intelligence isn't everything.
You saw how she acted today: she lacks the social conduct to lead other chipmunks."

``Maybe you should have spoken with her before announcing your decision to the whole Aorta?"

``It's not only today and you know it.
Never mind.
If not for us, then what of you and Cinnamon?"

``Don't," Kraerion said, flustered.

``What, you haven't told her?" Aequitius said, genuinely surprised.
He was just about to scold his son, but hesitated when he saw Kraerion's face and realized it would be superfluous.
Before he had time to say anything more, Kraerion turned around and ran.

Kraerion ran, and he ran, and he kept on running until the pain in his lungs and muscles numbed the thoughts that had so rudely occupied his mind.
